\chapter{Introduction}
\label{Chapter1}

Difference-in-Differences (DiD) has become one of the most widely used tools in applied microeconomics, largely because it offers an intuitive research design and can be implemented with \emph{two-way fixed effects} (TWFE) regressions \citep{imbens2009recent, athey2017state, currie2020technology, goldsmith2024tracking}. In settings with staggered adoption and treatment-effect heterogeneity, however, TWFE may fail to recover the causal parameter of interest and may instead combine comparisons in problematic ways, including through negative weights \citep{goodman2021difference, roth2023s}. Recent work has responded with a set of estimators designed for these settings, including \textcite{callaway2021difference, sun2021estimating, borusyak2024revisiting,dechaisemartin2020twfe,gardner2022two,wooldridge2025two}.

\textcite{baker2025difference} provide a general organizing framework that shifts attention away from familiar regressions and toward target parameters, identifying assumptions, and estimation strategies matched to the design. Their framework clarifies the structure of modern DiD designs, especially in settings with staggered treatment timing, alternative control groups, and nontrivial weighting. But much of the applied literature still absorbs this critique through a narrower lens. In practice, the discussion often collapses to negative weights alone, even though the empirical fragilities extend well beyond that margin.

This dissertation takes that gap between methodological guidance and applied practice as its starting point. In applied work, disagreement between TWFE and modern estimators is often shaped by a broader set of decisions, including the choice of control group, support at long horizons in event studies, dynamic contamination, the treatment of extreme horizons through \emph{binning}, the use of covariates and the move from unconditional to conditional parallel trends, the aggregation of dynamic effects, and implementation differences across software packages. The theoretical literature has clarified each of these issues. What remains less developed is how they appear jointly in published empirical work, how often they matter in practice, and which margins deserve priority when researchers read or implement a DiD design.

The dissertation addresses two questions. First, do conclusions from published DiD studies estimated with TWFE survive re-estimation with modern estimators? Second, which empirical margins account for the differences when they do not? I answer these questions through a large-scale reanalysis of 56 published articles that use DiD, have viable replication packages, and can be re-estimated with modern estimators. The analysis focuses on \emph{health economics}, where DiD has become increasingly common over the past decade \citep{goldsmith2024tracking}. The reanalysis uses AI support for screening, locating replication packages, and standardizing repetitive tasks, always under manual supervision. Chapter~\ref{Chapter3} discusses this workflow in detail and relates it to recent work on scalability in replication and reanalysis \parencite{xu2026scaling}.

The central result is that the literature does not collapse, but it often weakens and frequently shifts in ways that matter for interpretation. In roughly 86\% of cases, the sign of the effect is preserved under re-estimation, and nearly half also preserve statistical significance. At the same time, significance losses are common, magnitude changes are often large, and a substantial share of the differences come from margins other than negative weights alone. HonestDiD \citep{rambachan2023more} leads to a similar conclusion. Replacing TWFE with a heterogeneity-robust estimator may improve the interpretation of the estimand, but it does not rescue designs whose conclusions are fragile to modest departures from parallel trends. The practical implication is that modern estimators often change the strength of the evidence and sometimes its interpretation, even when they do not overturn the sign of the estimated effect.

The dissertation extends the reanalysis agenda developed in Finance and Accounting \parencite{baker2022much}, Political Science \parencite{chiu2023and}, and the broader reproducibility literature \parencite{brodeur2026reproducibility} to the \emph{health economics} literature that relies on DiD. It also shifts attention away from negative weights as the sole practical concern and toward a wider set of empirical margins that shape reanalysis results. The evidence from these re-estimations yields practical lessons for applied DiD work, not by proposing a new methodological framework, but by documenting where applied implementations most often drift away from the intended estimand, the intended comparison, or the available empirical support.

The rest of the dissertation proceeds as follows. Section~\ref{Chapter2} reviews the methodological and reanalysis literatures that motivate the exercise. Section~\ref{Chapter3} describes the sample construction, the reanalysis protocol, the metrics used to measure empirical change, and the role of AI support in screening, locating replication packages, and standardizing parts of the workflow. Section~\ref{Chapter4} presents the main results, organized around the core lessons of the reanalysis and the empirical mechanisms behind the observed changes. Section~\ref{Chapter5} concludes.
